\chapter{Graphs: Minimum Spanning Tree and Disjoint Set}
\chapterepigraph{A spanning tree connects every node using the fewest
possible edges, with no cycles. A minimum spanning tree does the same
at the lowest possible total weight. Answering ``would adding this
edge create a cycle'' efficiently -- without re-running a traversal
every time -- is what the Disjoint Set data structure is built for, and
once it exists, an entire family of ``group things by connectivity''
problems falls out almost for free.}

\section{The Disjoint Set Data Structure}
\label{sec:graph-disjoint-set}

A \textbf{Disjoint Set} (also called \textbf{Union-Find}) maintains a
collection of nodes partitioned into disjoint groups, supporting two
operations efficiently: $\textit{find}(x)$ (which group is $x$ in?)
and $\textit{union}(x,y)$ (merge $x$'s and $y$'s groups). Both run in
\textbf{nearly constant} amortized time with two optimizations:

\begin{ideabox}[Path Compression and Union by Size/Rank]
\begin{itemize}
  \item \textbf{Path compression}: during $\textit{find}(x)$, after
        walking up to the group's root, point every node visited along
        the way \textbf{directly} to the root -- flattening the tree
        so future lookups for those nodes are immediate.
  \item \textbf{Union by size (or rank)}: when merging two groups,
        attach the \textbf{smaller} group's root under the
        \textbf{larger} group's root, keeping the resulting tree
        shallow rather than letting it grow into a long chain.
\end{itemize}
Together, these two optimizations give nearly $O(1)$ amortized time
per operation (formally $O(\alpha(n))$, the inverse Ackermann
function, which is $\le 4$ for any practically sized input).
\end{ideabox}

\begin{patternbox}[What This Chapter Covers]
\begin{itemize}
  \item \textbf{Disjoint Set itself (Section~\ref{sec:graph-disjoint-set-impl})}.
  \item \textbf{Kruskal's (Section~\ref{sec:graph-kruskal})} and
        \textbf{Prim's (Section~\ref{sec:graph-prim})} algorithms for
        building a Minimum Spanning Tree -- one sorts edges and uses
        Disjoint Set to reject cycle-forming edges, the other grows a
        single tree greedily using a priority queue, mirroring
        Dijkstra's frontier expansion.
  \item \textbf{Number of Operations to Make Network Connected
        (Section~\ref{sec:graph-network-connected})}, \textbf{Most
        Stones Removed (Section~\ref{sec:graph-most-stones-removed})},
        \textbf{Accounts Merge (Section~\ref{sec:graph-accounts-merge})},
        \textbf{Number of Islands II (Section~\ref{sec:graph-number-of-islands-2})},
        and \textbf{Making a Large Island (Section~\ref{sec:graph-making-large-island})}:
        five problems that are not MST problems at all, but where
        Disjoint Set's efficient \emph{grouping} (not minimum-weight
        connecting) is exactly the tool needed.
\end{itemize}
\end{patternbox}

\newpage

\section{Disjoint Set (Union-Find)}
\label{sec:graph-disjoint-set-impl}

\problemstatement{Implement a data structure supporting: $\textit{find}(x)$
-- determine which group $x$ currently belongs to -- and
$\textit{union}(x,y)$ -- merge $x$'s and $y$'s groups into one, both
as efficiently as possible.}

\begin{patternbox}
\emph{``Efficiently track and merge groups of connected elements''} is
the textbook use case for Disjoint Set: each element points to a
\textbf{parent}, with a group's identity determined by walking parent
pointers up to a \textbf{root} (a node that is its own parent).
Path compression and union by size, described in
Section~\ref{sec:graph-disjoint-set}, keep both operations nearly
constant time.
\end{patternbox}

\subsection*{Algorithm}

\begin{enumerate}
  \item Initialize: every node is its own parent (its own group), and
        every group has size $1$.
  \item $\textit{find}(x)$: if $x$ is its own parent, return $x$.
        Otherwise, recursively find $x$'s ultimate root, and
        \textbf{set $x$'s parent directly to that root} before
        returning (path compression).
  \item $\textit{union}(x,y)$: find both roots; if equal, they are
        already in the same group, do nothing. Otherwise, attach the
        \textbf{smaller}-size root under the \textbf{larger}-size
        root's parent pointer, and add the smaller group's size into
        the larger's.
\end{enumerate}

\subsection*{C++ implementation}

\begin{lstlisting}
#include <bits/stdc++.h>
using namespace std;

class DisjointSet {
    vector<int> parent, size;
public:
    DisjointSet(int n) : parent(n), size(n, 1) {
        for (int i = 0; i < n; i++) parent[i] = i;
    }

    int find(int x) {
        if (parent[x] == x) return x;
        return parent[x] = find(parent[x]); // path compression
    }

    void unite(int x, int y) {
        int rootX = find(x), rootY = find(y);
        if (rootX == rootY) return;

        if (size[rootX] < size[rootY]) swap(rootX, rootY);
        parent[rootY] = rootX;       // smaller under larger
        size[rootX] += size[rootY];
    }

    bool connected(int x, int y) {
        return find(x) == find(y);
    }
};

int main() {
    DisjointSet ds(6);
    ds.unite(0, 1);
    ds.unite(1, 2);
    ds.unite(3, 4);

    cout << (ds.connected(0, 2) ? "true" : "false") << endl;  // true
    cout << (ds.connected(0, 3) ? "true" : "false") << endl;  // false
    ds.unite(2, 3);
    cout << (ds.connected(0, 4) ? "true" : "false") << endl;  // true
    return 0;
}
\end{lstlisting}

\subsection*{Dry run}

\dryrun{Take $6$ elements, performing $\textit{unite}(0,1)$,
$\textit{unite}(1,2)$, $\textit{unite}(3,4)$, then
$\textit{unite}(2,3)$.}

Initially each node is its own root, size $1$ each.
$\textit{unite}(0,1)$: both size $1$, $\textit{rootX}=0$ stays (tie,
no swap needed in this implementation since $1<1$ is false), parent$[1]=0$,
size$[0]=2$. $\textit{unite}(1,2)$: $\textit{find}(1)=0$ (size $2$),
$\textit{find}(2)=2$ (size $1$); $2$'s group is smaller, parent$[2]=0$,
size$[0]=3$. $\textit{unite}(3,4)$: similarly, parent$[4]=3$,
size$[3]=2$. $\textit{unite}(2,3)$: $\textit{find}(2)=0$ (size $3$,
via path compression directly), $\textit{find}(3)=3$ (size $2$); $3$'s
group is smaller, parent$[3]=0$, size$[0]=5$. Now $0,1,2,3,4$ are all
connected (root $0$); $5$ remains its own separate group.

\begin{complexitybox}
\textbf{Time Complexity.} $O(\alpha(n))$ amortized per
$\textit{find}$/$\textit{union}$ call, where $\alpha$ is the inverse
Ackermann function -- effectively constant for any realistic input
size.
\textbf{Space Complexity.} $O(n)$ for the parent and size arrays.
\end{complexitybox}

\begin{takeawaybox}
Disjoint Set answers ``are these two things in the same group, and can
I merge their groups'' far faster than re-running a DFS/BFS traversal
from scratch after every update -- making it the tool of choice
whenever a problem involves \textbf{repeated} connectivity queries or
incremental merging, rather than a single one-shot traversal.
\end{takeawaybox}

\section{Kruskal's Algorithm (Minimum Spanning Tree)}
\label{sec:graph-kruskal}

\problemstatement{Given a weighted, connected, undirected graph, find a
\textbf{Minimum Spanning Tree} (MST) -- a subset of edges connecting
all nodes, with no cycles, at the minimum possible total weight.}

\begin{patternbox}
\emph{``Connect everything as cheaply as possible''} is solved
greedily: sort all edges by weight, and consider them from
\textbf{cheapest to most expensive}, adding each edge to the MST
unless it would create a \textbf{cycle} -- checked instantly using
Section~\ref{sec:graph-disjoint-set-impl}'s Disjoint Set
($\textit{find}$ on both endpoints; if equal, they're already
connected, so this edge would close a cycle and is skipped).
\end{patternbox}

\subsection*{Why the greedy choice is correct (the cut property)}

For any partition of the nodes into two non-empty groups, the
\textbf{cheapest} edge crossing between them must belong to \emph{some}
MST -- if it didn't, swapping it in for whatever more expensive edge
\emph{does} cross between the groups in that MST would only decrease
the total weight, contradicting minimality. Processing edges from
cheapest to most expensive and greedily accepting any edge that
connects two still-separate components is exactly repeatedly applying
this cut property, which is why the greedy strategy provably produces
an MST.

\subsection*{Algorithm}

\begin{enumerate}
  \item Sort all edges by weight, ascending.
  \item Initialize a Disjoint Set over all $n$ nodes.
  \item For each edge $(u,v,w)$ in sorted order: if $\textit{find}(u)
        \neq \textit{find}(v)$, add the edge to the MST (accumulate
        its weight) and $\textit{unite}(u,v)$.
  \item Stop once $n{-}1$ edges have been added (a spanning tree on
        $n$ nodes always has exactly $n{-}1$ edges).
\end{enumerate}

\subsection*{C++ implementation}

\begin{lstlisting}
#include <bits/stdc++.h>
using namespace std;

class DisjointSet {
    vector<int> parent, size;
public:
    DisjointSet(int n) : parent(n), size(n, 1) {
        for (int i = 0; i < n; i++) parent[i] = i;
    }
    int find(int x) {
        if (parent[x] == x) return x;
        return parent[x] = find(parent[x]);
    }
    void unite(int x, int y) {
        int rootX = find(x), rootY = find(y);
        if (rootX == rootY) return;
        if (size[rootX] < size[rootY]) swap(rootX, rootY);
        parent[rootY] = rootX;
        size[rootX] += size[rootY];
    }
};

int kruskalMST(int n, vector<vector<int>>& edges) {
    sort(edges.begin(), edges.end(),
         [](const vector<int>& a, const vector<int>& b) { return a[2] < b[2]; });

    DisjointSet ds(n);
    int totalWeight = 0, edgesUsed = 0;

    for (auto& e : edges) {
        int u = e[0], v = e[1], w = e[2];
        if (ds.find(u) != ds.find(v)) {
            ds.unite(u, v);
            totalWeight += w;
            edgesUsed++;
            if (edgesUsed == n - 1) break;
        }
    }
    return totalWeight;
}

int main() {
    vector<vector<int>> edges = {{0,1,4},{0,2,1},{1,2,2},{1,3,5},{2,3,8}};
    cout << kruskalMST(4, edges) << endl; // 8
    return 0;
}
\end{lstlisting}

\subsection*{Dry run}

\dryrun{Take edges $0{-}1(4),0{-}2(1),1{-}2(2),1{-}3(5),2{-}3(8)$,
$n=4$.}

Sorted: $0{-}2(1),1{-}2(2),0{-}1(4),1{-}3(5),2{-}3(8)$. Take $0{-}2$
(weight $1$): different components, add; total $=1$. Take $1{-}2$
(weight $2$): $\textit{find}(1) \neq \textit{find}(2)$, add; total
$=3$; now $\{0,1,2\}$ connected. Take $0{-}1$ (weight $4$):
$\textit{find}(0)=\textit{find}(1)$ (both already in $\{0,1,2\}$'s
group) -- \textbf{skip}, would form a cycle. Take $1{-}3$ (weight
$5$): different components ($3$ is still separate), add; total
$=8$. Now $3$ edges used ($0{-}2$, $1{-}2$, $1{-}3$) $=n{-}1=3$, stop.
Final MST weight: $8$ -- node $3$'s cheapest available connection in
this graph is the weight-$5$ edge to node $1$, since the alternative
$2{-}3$ edge costs $8$.

\begin{complexitybox}
\textbf{Time Complexity.} $O(E \log E)$ for sorting, plus
$O(E \cdot \alpha(V))$ for the Disjoint Set operations:
$O(E \log E)$ overall.
\textbf{Space Complexity.} $O(V)$ for the Disjoint Set, plus $O(E)$ for
sorting the edge list.
\end{complexitybox}

\begin{takeawaybox}
Kruskal's algorithm is a direct, clean composition of two ideas already
established elsewhere in this book series: sorting to consider choices
in a greedy-optimal order (the same idea behind every Greedy chapter
algorithm), and Disjoint Set to reject cycle-forming choices in
near-constant time. Neither idea is new here -- their combination is.
\end{takeawaybox}

\section{Prim's Algorithm (Minimum Spanning Tree)}
\label{sec:graph-prim}

\problemstatement{Same as Section~\ref{sec:graph-kruskal}: find a
Minimum Spanning Tree of a weighted, connected, undirected graph --
this time by \textbf{growing a single tree} outward from one starting
node, rather than considering all edges globally sorted by weight.}

\begin{patternbox}
\emph{``Grow a tree, always extending it as cheaply as possible''}
mirrors the Graph Shortest Path chapter's Dijkstra's algorithm almost exactly: a
priority queue always offers the \textbf{cheapest available edge}
connecting some node already in the tree to some node not yet in it.
The key difference from Dijkstra: the priority queue tracks the
\textbf{edge weight to join the tree}, not the \textbf{cumulative
distance from a source} -- a node already close to the tree via one
cheap edge is preferred even if a different, longer path from the
start node exists.
\end{patternbox}

\subsection*{Why this also satisfies the cut property}

At every step, the tree built so far and the remaining nodes form a
cut (a partition into two groups). Prim's always picks the
\textbf{cheapest edge crossing that specific cut} -- exactly the
greedy choice the cut property (Section~\ref{sec:graph-kruskal})
guarantees is always safe to include in \emph{some} MST. Kruskal's
considers cuts implicitly through global edge sorting; Prim's
considers one specific, growing cut at each step.

\subsection*{Algorithm}

\begin{enumerate}
  \item Maintain a $\textit{inMST}$ boolean array and a min-heap of
        $(\textit{weight}, \textit{node})$ pairs, starting with
        $(0, \textit{anyStartNode})$.
  \item Pop the smallest-weight entry; if its node is already in the
        MST, skip (stale entry, exactly as in Dijkstra). Otherwise,
        mark it in the MST and add its weight to the total.
  \item For every neighbor of the newly added node not yet in the
        MST, push $(\textit{edgeWeight}, \textit{neighbor})$.
  \item Repeat until every node is in the MST.
\end{enumerate}

\subsection*{C++ implementation}

\begin{lstlisting}
#include <bits/stdc++.h>
using namespace std;

int primMST(int n, vector<vector<pair<int,int>>>& adj) {
    vector<bool> inMST(n, false);
    priority_queue<pair<int,int>, vector<pair<int,int>>, greater<>> pq;
    pq.push({0, 0}); // start from node 0, weight 0
    int totalWeight = 0;

    while (!pq.empty()) {
        auto [weight, node] = pq.top();
        pq.pop();
        if (inMST[node]) continue; // stale entry

        inMST[node] = true;
        totalWeight += weight;

        for (auto& [neighbor, edgeWeight] : adj[node]) {
            if (!inMST[neighbor]) {
                pq.push({edgeWeight, neighbor});
            }
        }
    }
    return totalWeight;
}

int main() {
    int n = 4;
    vector<vector<pair<int,int>>> adj(n);
    auto addEdge = [&](int u, int v, int w) {
        adj[u].push_back({v, w});
        adj[v].push_back({u, w});
    };
    addEdge(0, 1, 4); addEdge(0, 2, 1); addEdge(1, 2, 2);
    addEdge(1, 3, 5); addEdge(2, 3, 8);

    cout << primMST(n, adj) << endl; // 8
    return 0;
}
\end{lstlisting}

\subsection*{Dry run}

\dryrun{Same graph as Section~\ref{sec:graph-kruskal}'s example,
starting from node $0$.}

Heap: $[(0,0)]$. Pop $(0,0)$: add $0$ to MST, total $=0$; push
neighbors $(4,1)$ and $(1,2)$. Heap: $[(1,2),(4,1)]$. Pop $(1,2)$
(smallest): add $2$, total $=1$; push $2$'s neighbors not in MST:
$(2,1)$ (edge $1{-}2$, weight $2$) and $(8,3)$ (edge $2{-}3$, weight
$8$). Heap: $[(2,1),(4,1)_{\text{stale dup}},(8,3)]$. Pop $(2,1)$:
add $1$, total $=3$; push $1$'s remaining neighbor not in MST: $(5,3)$
(edge $1{-}3$, weight $5$). Heap now has both $(5,3)$ and $(8,3)$ for
node $3$. Pop $(4,1)_{\text{stale}}$: $1$ already in MST, skip. Pop
$(5,3)$ (smaller than the $(8,3)$ entry): add $3$, total $=8$. Final
MST weight: $8$ -- matching Kruskal's result on the same graph, as
expected (both algorithms always find \emph{a} minimum spanning tree,
and when the MST is unique, as here, the total weight necessarily
agrees).

\begin{complexitybox}
\textbf{Time Complexity.} $O(E \log V)$ -- each edge can trigger one
heap push, each costing $O(\log V)$.
\textbf{Space Complexity.} $O(V)$ for \textit{inMST}, plus $O(E)$ for
the heap in the worst case.
\end{complexitybox}

\begin{takeawaybox}
Kruskal's and Prim's solve the identical problem from two different
angles -- global edge sorting versus local frontier growth -- and
which is preferable depends on graph density: Kruskal's
$O(E\log E)$ favors sparse graphs (few edges to sort), while Prim's
$O(E\log V)$ with an adjacency-list-based heap favors graphs where $V$
is comparatively small relative to $E$. Both reuse machinery already
built earlier in this book: Kruskal's reuses Disjoint Set, Prim's
reuses Dijkstra's exact priority-queue frontier-expansion shape.
\end{takeawaybox}

\section{Number of Operations to Make Network Connected}
\label{sec:graph-network-connected}

\problemstatement{Given $n$ computers and a list of existing cable
connections, find the \textbf{minimum number of operations} to make
every computer connected, where one operation removes a cable and uses
it to connect any two currently-disconnected computers. Return $-1$ if
there are not enough cables to possibly connect everything.}

\begin{patternbox}
\emph{``Reuse redundant connections to join separate groups''} is a
direct Disjoint Set application: union every given connection, then
count how many separate \textbf{components} remain. Connecting $k$
components into $1$ always takes exactly $k{-}1$ operations (each
operation merges two components into one), and a \textbf{redundant}
cable (one connecting two computers \emph{already} in the same
component) is exactly a ``spare'' cable available to be moved.
\end{patternbox}

\subsection*{Why redundant cables exactly cover the requirement}

A graph with $n$ nodes and $c$ connected components needs exactly
$c{-}1$ more edges to become fully connected (this is the same
spanning-tree counting argument from Section~\ref{sec:graph-kruskal}:
each component is already internally connected, like a single
``super-node,'' and connecting $c$ super-nodes into one tree needs
$c{-}1$ edges). The number of redundant cables (edges that connected
two nodes already in the same component when processed) is exactly
$\textit{totalCables} - (n - c)$, since a spanning structure over $n$
nodes split into $c$ components uses exactly $n-c$ non-redundant
edges. So the operation is feasible exactly when the redundant count
is $\ge c{-}1$ -- equivalently, exactly when $\textit{totalCables}
\ge n-1$.

\subsection*{Algorithm}

\begin{enumerate}
  \item If the number of connections is less than $n{-}1$, return
        $-1$ immediately (not enough cable to possibly connect $n$
        computers).
  \item Union every connection using Disjoint Set.
  \item Count the number of distinct components $c$ remaining (count
        nodes that are their own root).
  \item Return $c{-}1$.
\end{enumerate}

\subsection*{C++ implementation}

\begin{lstlisting}
#include <bits/stdc++.h>
using namespace std;

class DisjointSet {
    vector<int> parent, size;
public:
    DisjointSet(int n) : parent(n), size(n, 1) {
        for (int i = 0; i < n; i++) parent[i] = i;
    }
    int find(int x) {
        if (parent[x] == x) return x;
        return parent[x] = find(parent[x]);
    }
    void unite(int x, int y) {
        int rootX = find(x), rootY = find(y);
        if (rootX == rootY) return;
        if (size[rootX] < size[rootY]) swap(rootX, rootY);
        parent[rootY] = rootX;
        size[rootX] += size[rootY];
    }
};

int makeConnected(int n, vector<vector<int>>& connections) {
    if ((int)connections.size() < n - 1) return -1;

    DisjointSet ds(n);
    for (auto& c : connections) ds.unite(c[0], c[1]);

    int components = 0;
    for (int i = 0; i < n; i++) {
        if (ds.find(i) == i) components++;
    }
    return components - 1;
}

int main() {
    vector<vector<int>> connections = {{0,1},{0,2},{1,2}};
    cout << makeConnected(4, connections) << endl; // 1
    return 0;
}
\end{lstlisting}

\subsection*{Dry run}

\dryrun{Take $n=4$, $\textit{connections}=[[0,1],[0,2],[1,2]]$.}

$3$ connections $\ge n{-}1=3$: feasible. Union $0{-}1$, $0{-}2$
(now $\{0,1,2\}$ connected), $1{-}2$ (redundant -- already same
component). After processing, components: $\{0,1,2\}$ (root, say,
$0$) and $\{3\}$ (its own root) -- $c=2$. Answer: $c-1=1$ -- one
redundant cable (the $1{-}2$ connection) can be moved to connect node
$3$.

\begin{complexitybox}
\textbf{Time Complexity.} $O((n+E)\cdot\alpha(n))$ -- unioning every
connection plus a final scan over all nodes.
\textbf{Space Complexity.} $O(n)$ for the Disjoint Set.
\end{complexitybox}

\begin{takeawaybox}
This is the first of several problems in this chapter where Disjoint
Set is used purely for its \textbf{grouping} ability -- counting and
merging components -- with no notion of edge weight or minimality
involved at all, distinguishing it clearly from the MST algorithms
earlier in this chapter that also relied on Disjoint Set.
\end{takeawaybox}

\section{Most Stones Removed with Same Row or Column}
\label{sec:graph-most-stones-removed}

\problemstatement{Given the coordinates of $n$ stones on a 2D plane, a
stone can be \textbf{removed} if some \textbf{other}, not-yet-removed
stone shares its row or its column. Find the maximum number of stones
that can be removed.}

\begin{patternbox}
\emph{``Stones connected by shared rows/columns''} is a Disjoint Set
grouping problem in disguise: union every pair of stones sharing a row
or column, then within each resulting connected group, \textbf{every
stone but one} can always be removed (remove them one at a time,
always leaving at least one survivor in the group to justify removing
the next). The answer is simply
$\textit{totalStones} - \textit{numberOfGroups}$.
\end{patternbox}

\subsection*{Why one survivor per group is always achievable}

Within a connected group of stones (connected via a chain of shared
rows/columns), pick any spanning tree of that connectivity (exactly
as in Section~\ref{sec:graph-kruskal}'s spanning tree idea) and remove
stones as \textbf{leaves} of that tree, working inward: a leaf stone
always shares its row or column with its tree-parent, which is
guaranteed not yet removed (it's deeper in the removal order), so the
removal condition is always satisfiable until only the root remains.
This is why the answer doesn't depend on \emph{how} stones happen to
connect within a group -- only on \emph{how many} groups there are.

\subsection*{Building the Disjoint Set over rows and columns}

Rather than unioning stones pairwise directly (which would need
checking every pair, $O(n^2)$), union each stone's \textbf{row} and
\textbf{column} together in a combined coordinate space: treat row
indices and column indices as Disjoint Set elements in two disjoint
ranges (e.g., column $c$ mapped to $c + 10001$ to avoid colliding with
row indices). Two stones sharing a row or column end up in the same
Disjoint Set group automatically, transitively, without ever comparing
stones to each other directly.

\subsection*{C++ implementation}

\begin{lstlisting}
#include <bits/stdc++.h>
using namespace std;

class DisjointSet {
    unordered_map<int,int> parent;
public:
    int find(int x) {
        if (parent.find(x) == parent.end()) parent[x] = x;
        if (parent[x] == x) return x;
        return parent[x] = find(parent[x]);
    }
    void unite(int x, int y) {
        parent[find(x)] = find(y);
    }
};

int removeStones(vector<vector<int>>& stones) {
    DisjointSet ds;
    for (auto& s : stones) {
        ds.unite(s[0], s[1] + 10001); // column shifted into a separate range
    }

    unordered_set<int> roots;
    for (auto& s : stones) {
        roots.insert(ds.find(s[0]));
    }
    return (int)stones.size() - (int)roots.size();
}

int main() {
    vector<vector<int>> stones = {{0,0},{0,1},{1,0},{1,2},{2,1},{2,2}};
    cout << removeStones(stones) << endl; // 5
    return 0;
}
\end{lstlisting}

\subsection*{Dry run}

\dryrun{Take $\textit{stones}=[(0,0),(0,1),(1,0),(1,2),(2,1),(2,2)]$.}

Union (row $0$, col $0{+}10001$), (row $0$, col $1{+}10001$), (row
$1$, col $0{+}10001$), (row $1$, col $2{+}10001$), (row $2$, col
$1{+}10001$), (row $2$, col $2{+}10001$). Tracing connectivity: stone
$(0,0)$ links row $0$ and col $10001$; stone $(1,0)$ links row $1$ and
col $10001$ -- so row $0$ and row $1$ end up in the \emph{same} group
via shared column $0$. Continuing, every row and column index here
ends up merged into a \textbf{single} group (rows $0,1,2$ and columns
$0,1,2$ all connected through the chain of shared coordinates). All
$6$ stones' row-roots resolve to this one group: $\textit{roots}$ has
size $1$. Answer: $6 - 1 = 5$.

\begin{complexitybox}
\textbf{Time Complexity.} $O(n \cdot \alpha(n))$ -- unioning and
finding over $n$ stones.
\textbf{Space Complexity.} $O(n)$ for the Disjoint Set's hash-map-based
parent storage (rows and columns are sparse, so a hash map avoids
allocating a dense array over the full coordinate range).
\end{complexitybox}

\begin{takeawaybox}
Mapping two conceptually different categories (rows and columns) into
a single Disjoint Set by shifting one category's indices into a
disjoint numeric range is a general technique whenever ``connected
through a shared attribute'' needs to be modeled without explicitly
enumerating every pairwise connection.
\end{takeawaybox}

\section{Accounts Merge}
\label{sec:graph-accounts-merge}

\problemstatement{Given a list of accounts, each formatted as
$[\textit{name}, \textit{email}_1, \textit{email}_2, \ldots]$, merge
any accounts that share \textbf{at least one} email in common (the
same person may have multiple accounts with overlapping email sets).
Return the merged accounts, each with its emails sorted, name listed
first.}

\begin{patternbox}
\emph{``Merge groups that share any common element''} is Disjoint Set
over \textbf{account indices}: for every email seen, if it has already
been seen belonging to some earlier account, union the current account
with that earlier one. Accounts ending up with the same root after all
unions belong to the same person and should be merged.
\end{patternbox}

\subsection*{Algorithm}

\begin{enumerate}
  \item Initialize a Disjoint Set over account \textbf{indices}
        (\emph{not} emails -- emails are not the nodes here, accounts
        are).
  \item Maintain a hash map from email to the \textbf{first} account
        index seen owning it. For each account $i$, for each of its
        emails: if the email has been seen before (owned by account
        $j$), $\textit{unite}(i,j)$; otherwise, record $i$ as this
        email's owner.
  \item Group every email by the \textbf{root} of its originating
        account index.
  \item For each group, output $[\textit{name}, \text{sorted emails}]$
        (the name can be taken from any account in the group, since
        merged accounts are assumed to share the same name).
\end{enumerate}

\subsection*{C++ implementation}

\begin{lstlisting}
#include <bits/stdc++.h>
using namespace std;

class DisjointSet {
    vector<int> parent;
public:
    DisjointSet(int n) : parent(n) {
        for (int i = 0; i < n; i++) parent[i] = i;
    }
    int find(int x) {
        if (parent[x] == x) return x;
        return parent[x] = find(parent[x]);
    }
    void unite(int x, int y) {
        parent[find(x)] = find(y);
    }
};

vector<vector<string>> accountsMerge(vector<vector<string>>& accounts) {
    int n = accounts.size();
    DisjointSet ds(n);
    unordered_map<string, int> emailOwner;

    for (int i = 0; i < n; i++) {
        for (int j = 1; j < (int)accounts[i].size(); j++) {
            string& email = accounts[i][j];
            if (emailOwner.find(email) != emailOwner.end()) {
                ds.unite(i, emailOwner[email]);
            } else {
                emailOwner[email] = i;
            }
        }
    }

    unordered_map<int, set<string>> grouped;
    for (auto& [email, ownerIndex] : emailOwner) {
        grouped[ds.find(ownerIndex)].insert(email);
    }

    vector<vector<string>> result;
    for (auto& [root, emails] : grouped) {
        vector<string> merged;
        merged.push_back(accounts[root][0]); // name
        merged.insert(merged.end(), emails.begin(), emails.end());
        result.push_back(merged);
    }
    return result;
}

int main() {
    vector<vector<string>> accounts = {
        {"John", "johnsmith@mail.com", "john_newyork@mail.com"},
        {"John", "johnsmith@mail.com", "john00@mail.com"},
        {"Mary", "mary@mail.com"}
    };
    for (auto& acc : accountsMerge(accounts)) {
        for (string& s : acc) cout << s << " ";
        cout << endl;
    }
    return 0;
}
\end{lstlisting}

\subsection*{Dry run}

\dryrun{Take the three accounts above.}

Account $0$: email \texttt{johnsmith@mail.com} unseen, owner $=0$;
email \texttt{john\_newyork@mail.com} unseen, owner $=0$. Account $1$:
email \texttt{johnsmith@mail.com} \textbf{already owned by account
$0$} -- $\textit{unite}(1,0)$; email \texttt{john00@mail.com} unseen,
owner $=1$. Account $2$: email \texttt{mary@mail.com} unseen, owner
$=2$. Grouping by root: accounts $0$ and $1$ share a root (merged via
the shared email), contributing all three of their emails sorted
together under name \texttt{"John"}; account $2$ stands alone under
\texttt{"Mary"}. Final: two merged account groups.

\begin{complexitybox}
\textbf{Time Complexity.} $O(\sum |\textit{emails}| \cdot \alpha(n) +
\sum |\textit{emails}| \log(\sum |\textit{emails}|))$ -- unioning
costs near-constant time per email, and sorting each group's emails
(via \texttt{set}) contributes the logarithmic factor.
\textbf{Space Complexity.} $O(\sum |\textit{emails}|)$ for the email
ownership map and grouped output.
\end{complexitybox}

\begin{takeawaybox}
The subtlety in this problem is recognizing that the Disjoint Set's
\textbf{nodes} are accounts, not emails -- emails are merely the
\emph{evidence} used to decide which accounts to union. This same
\textit{``the connectivity evidence and the thing being grouped are
different''} pattern reappears throughout real-world entity-resolution
problems (matching user profiles, deduplicating records) well beyond
this specific email-based example.
\end{takeawaybox}

\section{Number of Islands II}
\label{sec:graph-number-of-islands-2}

\problemstatement{Given an $m \times n$ grid initially entirely water,
and a sequence of positions where land is added \textbf{one cell at a
time}, return the number of islands \textbf{after each addition}.}

\begin{patternbox}
\emph{``Track connected components as elements are added one at a
time''} is exactly what Disjoint Set is built for, in contrast to the
Graph Basics chapter's Number of Islands (a single, static, one-shot
count via traversal): here, the grid \textbf{changes incrementally},
and re-running a full DFS/BFS after every single addition would cost
$O(mn)$ per step. Disjoint Set instead updates the island count in
near-constant time per addition.
\end{patternbox}

\subsection*{Algorithm}

\begin{enumerate}
  \item Maintain a Disjoint Set over all grid cells (indexed as
        $r \times \textit{cols} + c$), and an $\textit{islandCount}$
        starting at $0$.
  \item For each new land position $(r,c)$: if it's already land
        (a duplicate addition), record the current count unchanged and
        continue. Otherwise, mark it land and increment
        $\textit{islandCount}$ by $1$ (a brand new, isolated island).
  \item Check all $4$ neighbors: for each neighbor that is already
        land and in a \textbf{different} Disjoint Set group, union
        them and \textbf{decrement} $\textit{islandCount}$ by $1$
        (two previously separate islands just merged into one).
  \item Record $\textit{islandCount}$ as the answer for this step.
\end{enumerate}

\subsection*{C++ implementation}

\begin{lstlisting}
#include <bits/stdc++.h>
using namespace std;

class DisjointSet {
    vector<int> parent, size;
public:
    DisjointSet(int n) : parent(n), size(n, 1) {
        for (int i = 0; i < n; i++) parent[i] = i;
    }
    int find(int x) {
        if (parent[x] == x) return x;
        return parent[x] = find(parent[x]);
    }
    void unite(int x, int y) {
        int rootX = find(x), rootY = find(y);
        if (rootX == rootY) return;
        if (size[rootX] < size[rootY]) swap(rootX, rootY);
        parent[rootY] = rootX;
        size[rootX] += size[rootY];
    }
};

vector<int> numIslands2(int m, int n, vector<vector<int>>& positions) {
    DisjointSet ds(m * n);
    vector<vector<bool>> isLand(m, vector<bool>(n, false));
    int islandCount = 0;
    vector<int> result;

    int dr[] = {-1, 1, 0, 0};
    int dc[] = {0, 0, -1, 1};

    for (auto& pos : positions) {
        int r = pos[0], c = pos[1];
        if (isLand[r][c]) {
            result.push_back(islandCount);
            continue;
        }

        isLand[r][c] = true;
        islandCount++;

        for (int dir = 0; dir < 4; dir++) {
            int nr = r + dr[dir], nc = c + dc[dir];
            if (nr >= 0 && nr < m && nc >= 0 && nc < n && isLand[nr][nc]) {
                int idA = r * n + c, idB = nr * n + nc;
                if (ds.find(idA) != ds.find(idB)) {
                    ds.unite(idA, idB);
                    islandCount--;
                }
            }
        }
        result.push_back(islandCount);
    }
    return result;
}

int main() {
    vector<vector<int>> positions = {{0,0},{0,1},{1,2},{1,0},{0,2}};
    for (int x : numIslands2(3, 3, positions)) cout << x << " ";
    cout << endl; // 1 1 2 2 1
    return 0;
}
\end{lstlisting}

\subsection*{Dry run}

\dryrun{Take $m=n=3$,
$\textit{positions}=[(0,0),(0,1),(1,2),(1,0),(0,2)]$.}

Add $(0,0)$: new island, count $=1$. No land neighbors. Result: $1$.
Add $(0,1)$: new island, count $=2$; neighbor $(0,0)$ is land, union,
count $=1$. Result: $1$. Add $(1,2)$: new island, count $=2$; no land
neighbors yet. Result: $2$. Add $(1,0)$: new island, count $=3$;
neighbor $(0,0)$ is land (same group as $(0,1)$), union, count $=2$.
Result: $2$. Add $(0,2)$: new island, count $=3$; neighbor $(1,2)$ is
land, union, count $=2$; neighbor $(0,1)$ is land and -- after the
previous union -- in a \textbf{different} group from $(0,2)$'s newly
merged group, union again, count $=1$. Result: $1$. Final sequence:
$1,1,2,2,1$ -- the last addition, $(0,2)$, happens to bridge
\emph{two} previously separate islands at once, dropping the count by
$2$ in a single step.

\begin{complexitybox}
\textbf{Time Complexity.} $O(k \cdot \alpha(mn))$ for $k$ position
updates, each doing a constant number of Disjoint Set operations.
\textbf{Space Complexity.} $O(mn)$ for the Disjoint Set and the
\textit{isLand} grid.
\end{complexitybox}

\begin{takeawaybox}
This problem is the clearest illustration in this chapter of Disjoint
Set's core advantage over re-traversal: when connectivity needs to be
queried or updated \textbf{repeatedly} as a structure changes
incrementally, Disjoint Set answers each update in near-constant time,
where a fresh DFS/BFS per update would be asymptotically far slower.
\end{takeawaybox}

\section{Making a Large Island}
\label{sec:graph-making-large-island}

\problemstatement{Given an $n \times n$ binary grid, change \textbf{at
most one} cell from $0$ to $1$ to maximize the size of the largest
resulting island ($4$-connected group of $1$s). Return that maximum
possible size.}

\begin{patternbox}
This closing problem of the chapter combines Disjoint Set's component
\textbf{sizing} (already used implicitly by union-by-size throughout
this chapter) with a final scan: first, group every existing island
using Disjoint Set, tracking each group's size; then, for every water
cell, compute what size would result from flipping \textbf{just that
one cell} -- the sum of the sizes of its (up to $4$) \textbf{distinct}
neighboring islands, plus $1$ for the flipped cell itself.
\end{patternbox}

\subsection*{Why distinctness of neighboring islands matters}

A water cell might have two or more $4$-directional neighbors that
belong to the \emph{same} island (e.g., a water cell with land both
above and below, where that land happens to already be connected via
some other path around the grid). Counting that island's size twice
would overcount. So neighboring islands must be collected into a
\textbf{set} of distinct roots before summing their sizes -- exactly
the kind of deduplication Disjoint Set's $\textit{find}$ operation
makes trivial (two neighbors are ``the same island'' precisely when
$\textit{find}$ returns the same root for both).

\subsection*{Algorithm}

\begin{enumerate}
  \item Build a Disjoint Set over all $n^2$ cells; union every land
        cell with its land neighbors.
  \item For every water cell, collect the distinct roots among its up
        to $4$ neighbors that are land; the candidate size if flipped
        is $1$ plus the sum of those distinct roots' group sizes.
  \item Track the maximum candidate size across all water cells (and
        separately, the largest existing island's size, in case the
        grid has no water cell to flip at all -- the answer is then
        simply $n^2$).
\end{enumerate}

\subsection*{C++ implementation}

\begin{lstlisting}
#include <bits/stdc++.h>
using namespace std;

class DisjointSet {
    vector<int> parent, size;
public:
    DisjointSet(int n) : parent(n), size(n, 1) {
        for (int i = 0; i < n; i++) parent[i] = i;
    }
    int find(int x) {
        if (parent[x] == x) return x;
        return parent[x] = find(parent[x]);
    }
    void unite(int x, int y) {
        int rootX = find(x), rootY = find(y);
        if (rootX == rootY) return;
        if (size[rootX] < size[rootY]) swap(rootX, rootY);
        parent[rootY] = rootX;
        size[rootX] += size[rootY];
    }
    int getSize(int x) { return size[find(x)]; }
};

int largestIsland(vector<vector<int>>& grid) {
    int n = grid.size();
    DisjointSet ds(n * n);
    int dr[] = {-1, 1, 0, 0};
    int dc[] = {0, 0, -1, 1};

    for (int r = 0; r < n; r++) {
        for (int c = 0; c < n; c++) {
            if (grid[r][c] == 1) {
                for (int dir = 0; dir < 4; dir++) {
                    int nr = r + dr[dir], nc = c + dc[dir];
                    if (nr >= 0 && nr < n && nc >= 0 && nc < n && grid[nr][nc] == 1) {
                        ds.unite(r * n + c, nr * n + nc);
                    }
                }
            }
        }
    }

    int best = 0;
    bool hasWater = false;

    for (int r = 0; r < n; r++) {
        for (int c = 0; c < n; c++) {
            if (grid[r][c] == 0) {
                hasWater = true;
                unordered_set<int> distinctRoots;
                for (int dir = 0; dir < 4; dir++) {
                    int nr = r + dr[dir], nc = c + dc[dir];
                    if (nr >= 0 && nr < n && nc >= 0 && nc < n && grid[nr][nc] == 1) {
                        distinctRoots.insert(ds.find(nr * n + nc));
                    }
                }
                int candidate = 1;
                for (int root : distinctRoots) candidate += ds.getSize(root);
                best = max(best, candidate);
            } else {
                best = max(best, ds.getSize(r * n + c));
            }
        }
    }
    return hasWater ? best : n * n;
}

int main() {
    vector<vector<int>> grid = {
        {1, 0},
        {0, 1}
    };
    cout << largestIsland(grid) << endl; // 3
    return 0;
}
\end{lstlisting}

\subsection*{Dry run}

\dryrun{Take the $2\times2$ grid $[[1,0],[0,1]]$ (two separate
single-cell islands, diagonally placed).}

After building the Disjoint Set: cell $(0,0)$ and cell $(1,1)$ are
each their own island of size $1$ (they are not $4$-directionally
adjacent to each other). Checking water cell $(0,1)$: neighbors
$(0,0)$ (land, size $1$) and $(1,1)$ (land, size $1$) -- two
\emph{distinct} roots, candidate $=1+1+1=3$. Checking water cell
$(1,0)$: neighbors $(0,0)$ (land, size $1$) and $(1,1)$ (land, size
$1$) -- also distinct, candidate $=1+1+1=3$. Best: $3$ -- flipping
either water cell connects both existing islands into one.

\begin{complexitybox}
\textbf{Time Complexity.} $O(n^2 \cdot \alpha(n^2))$ -- building the
Disjoint Set and the final scan both touch every cell a constant
number of times.
\textbf{Space Complexity.} $O(n^2)$ for the Disjoint Set.
\end{complexitybox}

\begin{takeawaybox}
This chapter closes by combining nearly everything it introduced:
Disjoint Set's grouping (Section~\ref{sec:graph-disjoint-set-impl}),
size tracking (the same union-by-size bookkeeping Kruskal's and Prim's
relied on internally), and the deduplication that distinguishes a
\emph{group} from a list of individual elements (echoing Accounts
Merge's account-grouping and Most Stones Removed's row/column
grouping) -- a fitting capstone for a chapter built entirely around one
data structure's surprising range of applications.
\end{takeawaybox}

